% KHÔNG \input file này: attention / encoder-decoder dễ trùng bài viết Vaswani
% và blog mô hình. Chương 2 đã nói cửa sổ ngữ cảnh và nhiệt độ bằng ví dụ kho luật.

\section{Transformer, attention và mô hình nhúng --- mức đủ dùng}

Đồ án không huấn luyện Transformer. Mục này giải thích vừa đủ để lý giải các
tham số đã chọn (cửa sổ ngữ cảnh, nhiệt độ, số chiều embedding).

\subsection{Cơ chế attention}

Với một chuỗi token, attention tính trọng số: token hiện tại ``nhìn'' các vị trí
khác nhiều hay ít. Self-attention trong một lớp cho phép quan hệ xa (chủ ngữ
đầu câu với động từ cuối câu) tốt hơn mạng hồi quy thuần. Đa đầu (multi-head)
học nhiều kiểu quan hệ song song. Stack nhiều lớp tạo mô hình sâu~\cite{vaswani2017}.

Hệ quả thực tiễn: chi phí attention tăng theo bình phương độ dài chuỗi. Vì vậy
không nhét \soDieu{} Điều vào một prompt. Pipeline đồ án cắt còn $k=8$ sau lọc.

\subsection{Mô hình mã hoá và mô hình giải mã}

Encoder (BERT-style) đọc hai chiều, hợp câu phân loại và nhúng. Decoder
(GPT-style) sinh token trái sang phải, hợp hội thoại. Embedding API thường là
encoder hoặc biến thể. Gemini flash dùng khi sinh câu là decoder. Trộn nhầm
``một mô hình làm cả hai cổng'' dễ dẫn tới gọi sai endpoint: đồ án tách
\texttt{LLM\_MODEL} và \texttt{EMBEDDINGS\_MODEL}.

\subsection{Học biểu diễn phân tán}

Word2Vec gán mỗi từ một vector tĩnh, không phân biệt ``đóng'' (bảo hiểm) và
``đóng'' (cửa). Mô hình ngữ cảnh (contextual) gán vector theo câu. Embedding
Điều của đồ án là vector cả đoạn, không phải trung bình word2vec. Đó là lý do
cần API nhúng hiện đại, không tự huấn luyện skip-gram trên 38 luật (quá ít dữ
liệu, không bắt được ``sa thải'' $\leftrightarrow$ ``đơn phương chấm dứt'' nếu
cặp không đồng xuất hiện đủ).

\subsection{Nhiệt độ, top-p, số token tối đa}

Nhiệt độ nhân vào logits trước softmax. Top-p (nucleus) cắt đuôi phân bố. Đồ án
giữ nhiệt độ 0, không chỉnh top-p trên sản phẩm, ưu tiên ổn định diễn giải.
\texttt{max\_tokens} 4096 đủ một câu tư vấn kèm gạch đầu dòng, không đủ để mô
hình chép cả chương luật --- vốn không mong muốn.

\subsection{Hallucination dưới góc xác suất}

Mô hình tối ưu likelihood token, không tối ưu ``đúng Điều 44''. Khi ngữ cảnh
mỏng, phân bố vẫn gán xác suất cao cho các số hiệu hay gặp trong dữ liệu huấn
luyện. RAG tăng xác suất các token nằm trong đoạn đã nhét, nhưng không đưa về
không. Hậu kiểm citation là lớp xác định (deterministic) bù cho lớp xác suất.

\section{Truy hồi lai: vì sao không chỉ một kênh}

\subsection{Thất bại điển hình của từng kênh}

BM25 thất bại khi từ vựng lệch hoàn toàn (sa thải / đơn phương). Dense thất bại
khi truy vấn là số hiệu hoặc cụm cực ngắn, và khi hai Điều kề nhau nội dung na
ná (cùng chương thuế). RRF không sửa nội dung Điều; RRF chỉ tăng xác suất Điều
đúng có mặt trong union hai danh sách.

\subsection{Tham số $k$ trước và sau gộp}

Lấy $k_{\mathrm{RRF}}=60$ rồi cắt 8 là phễu: rộng lúc hợp nhất, hẹp lúc nhét
prompt. Quá hẹp lúc BM25/Chroma thì RRF không có gì để gộp. Quá rộng lúc prompt
thì filter và generate đắt, mô hình lạc. Con số 8 là thỏa hiệp cửa sổ và quan
sát tay, không phải cực trị trên tập nhãn.

\subsection{Cache BM25}

Tách từ \soDieu{} Điều tốn thời gian khởi động. Cache theo vân tay tokenizer:
đổi underthesea hoặc corpus thì tính lại. Không cache Chroma cùng file này vì
Chroma đã là chỉ mục trên đĩa volume.

\section{Kiến trúc REST, tài nguyên và idempotency}

Idempotent: gọi lại cùng request không đổi kết quả ngoài lần đầu (GET, PUT,
DELETE lý tưởng). POST generate lịch được viết idempotent theo (rule, kỳ, tổ
chức): đây là POST có ngữ nghĩa ``đảm bảo tồn tại'', gần PUT. Chat POST không
idempotent: mỗi lần một tin nhắn mới. Phân biệt này giải thích vì sao retry
client được phép trên generate lịch mà phải hỏi người dùng trước khi gửi lại
câu chat (tránh nhân đôi câu hỏi --- tuy câu hỏi đã persist theo id).
